\chapter{Zadatak 2 --- svođenje na konačno mnogo nejednakosti}
\label{ch:z2-vrhovi}

\section{Što se ovdje radi}

Pokazuje se da je supremum linearne funkcije po pravokutniku postignut u nekom
od četiriju vrhova, pa se beskonačan skup ograničenja svodi na osam običnih
linearnih nejednakosti. Ovo je matematička srž Zadatka 2.

\begin{predznanje}
\begin{itemize}[nosep,leftmargin=*]
  \item Skupovi $\mathcal{P}_1$, $\mathcal{P}_2$ --- \eqref{eq:z2-P}.
  \item Uvjet sa supremumom --- \eqref{eq:z2-sup}.
  \item Konveksna kombinacija i konveksna ljuska --- slajd niže.
\end{itemize}
\end{predznanje}

\section{Konveksna ljuska}
\label{sec:z2-ljuska}

\begin{izvor}{Predavanje 2 --- Definicije, klasifikacija, konveksnost, slajd 59/71 (PDF str.~60)}
\slajd[0.86]{p02-konveksna-ljuska}

\textbf{Transkripcija.} \emph{Konveksna kombinacija} točaka (vektora)
$x_1, x_2, \dots, x_n$ je točka (vektor) $x$ u formi
\[
x = \gamma_1 x_1 + \gamma_2 x_2 + \dots + \gamma_n x_n,
\qquad \gamma_1 + \dots + \gamma_n = 1, \quad \gamma_i \geq 0 .
\]
\emph{Konveksna ljuska} skupa $S$ je skup svih konveksnih kombinacija točaka iz
$S$. Konveksna ljuska skupa $S$ je najmanji konveksni skup koji sadrži $S$. Skup
je konveksan ako i samo ako je jednak svojoj konveksnoj ljusci.

\textbf{Zašto je to ovdje ključno:} pravokutnik je konveksan, pa je jednak
konveksnoj ljusci svojih četiriju vrhova. Svaka točka pravokutnika, dakle svaka
moguća realizacija koeficijenata, može se zapisati kao težinski prosjek vrhova s
nenegativnim težinama koje se zbrajaju u jedan.
\end{izvor}

\section{Teorem o vrhovima}
\label{sec:z2-teorem-vrhovi}

\begin{formula}{Supremum linearne funkcije po poliedru}
Neka je $\mathcal{P} = \mathrm{co}\bigl(v^{(1)}, \dots, v^{(N)}\bigr)$ konveksna
ljuska konačno mnogo točaka. Tada za svaki $x$ vrijedi
\begin{equation}
\sup_{a \in \mathcal{P}} a\trans x \leq h
\qquad\Longleftrightarrow\qquad
\bigl(v^{(k)}\bigr)\trans x \leq h \quad \text{za } k = 1, \dots, N .
\label{eq:z2-teorem-vrhovi}
\end{equation}
\end{formula}

\begin{korak}[Dokaz, oba smjera]
\begin{enumerate}[leftmargin=*]
  \item \textbf{($\Leftarrow$) Iz vrhova slijedi cijeli skup.} Neka je
        $a \in \mathcal{P}$ proizvoljan. Po definiciji konveksne ljuske postoje
        težine $\gamma_k \geq 0$, $\sum_k \gamma_k = 1$, takve da je
        $a = \sum_k \gamma_k v^{(k)}$. Tada je, zbog linearnosti skalarnog
        produkta po prvom argumentu,
        \[
          a\trans x = \sum_k \gamma_k \bigl(v^{(k)}\bigr)\trans x
                   \;\leq\; \sum_k \gamma_k\, h
                   \;=\; h \sum_k \gamma_k \;=\; h .
        \]
        Ključan je korak u sredini: svaki član je najviše $h$ po pretpostavci, a
        težine su nenegativne, pa se nejednakost smije množiti njima i zbrajati.
  \item \textbf{($\Rightarrow$) Iz cijelog skupa slijede vrhovi.} Trivijalno:
        vrhovi \emph{jesu} elementi skupa $\mathcal{P}$, pa uvjet koji vrijedi za
        sve članove vrijedi napose i za njih.
\end{enumerate}
Ekvivalencija je \textbf{egzaktna} --- nije riječ o aproksimaciji ni o
konzervativnoj ocjeni. Robusni problem i problem s vrhovima imaju identičan
dopustivi skup.
\end{korak}

\begin{zamka}[Zamka: ,,uzeo sam najgori slučaj, valjda je to dovoljno``]
Argument mora ići preko konveksne kombinacije. Nije dovoljno reći ,,linearna
funkcija je monotona pa je najgore u kutu`` --- to je intuicija, ne dokaz, i pada
čim skup nesigurnosti nije pravokutnik. Za elipsoidnu nesigurnost, na primjer,
vrhova nema, a supremum se računa preko norme; Predavanje 8 to obrađuje na
zadnja dva slajda.
\end{zamka}

\section{Osam nejednakosti}
\label{sec:z2-osam}

\begin{korak}[Nabrajanje vrhova]
Svaki pravokutnik u ravnini ima $2^2 = 4$ vrha --- sve kombinacije donje i gornje
granice po objema koordinatama. Uz dva ograničenja to daje $2 \cdot 4 = 8$
linearnih nejednakosti:
\begin{equation}
\begin{aligned}
-0{,}85 x_1 - 0{,}85 x_2 &\leq -10, &\qquad 0{,}85 x_1 - 0{,}85 x_2 &\leq 3, \\
-0{,}85 x_1 - 1{,}15 x_2 &\leq -10, &\qquad 0{,}85 x_1 - 1{,}15 x_2 &\leq 3, \\
-1{,}15 x_1 - 0{,}85 x_2 &\leq -10, &\qquad 1{,}15 x_1 - 0{,}85 x_2 &\leq 3, \\
-1{,}15 x_1 - 1{,}15 x_2 &\leq -10, &\qquad 1{,}15 x_1 - 1{,}15 x_2 &\leq 3.
\end{aligned}
\label{eq:z2-osam}
\end{equation}
\end{korak}

\begin{formula}{Robusni problem kao obični QP}
\begin{equation}
\min_x \ \tfrac12 x\trans Q x + c\trans x
\quad \text{uz} \quad
G_{\mathrm{rob}}\, x \leq h_{\mathrm{rob}},
\label{eq:z2-qp}
\end{equation}
gdje je $G_{\mathrm{rob}} \in \R^{8\times2}$ sastavljena od vrhova, a
$h_{\mathrm{rob}} = [-10,\,-10,\,-10,\,-10,\ 3,\,3,\,3,\,3]\trans$. Funkcija
cilja je nepromijenjena, pa je ovo \textbf{isti tip problema kao Zadatak 1} ---
samo s osam ograničenja umjesto dva.
\end{formula}

\begin{izvor}{Predavanje 8 --- Robusno optimiranje, slajd 16/21 (PDF str.~16)}
\slajd[0.86]{p08-primjer2-konacno}

\textbf{Što je na slajdu.} Profesorov Primjer 2: robusni problem prepisan kao
\emph{ekvivalentan LP problem (konačan broj nejednakosti!)}. Nominalna matrica
$A$ zbraja se s perturbacijskim matricama u kojima varijacije $\delta_i$ dolaze
na granice svojih intervala ($\pm 0{,}05$).

\textbf{Podudarnost s našim postupkom:} isti korak, samo drukčije zapisan.
Profesor perturbacije stavlja na granice intervala; mi nabrajamo vrhove
pravokutnika. To je isto: vrh pravokutnika \emph{jest} kombinacija granica.
Uskličnik u naslovu slajda je poanta cijelog zadatka --- beskonačno je postalo
konačno.
\end{izvor}

\section{Profesorova općenita metoda i zašto je ovdje ne trebamo}
\label{sec:z2-dualna-metoda}

Ovo je najvjerojatnije potpitanje na obrani, jer Predavanje 8 glavni rezultat
izvodi \emph{drugim} putem od onoga kojim je zadatak riješen.

\begin{izvor}{Predavanje 8 --- Robusno optimiranje, slajd 3/21 (PDF str.~3)}
\slajd[0.86]{p08-ekvivalentan-lp}

\textbf{Transkripcija.} Robusan LP
$\min_x c\trans x$ uz $a_i\trans x \leq b_i$ za sve
$a_i \in \mathcal{P}_i = \{a \mid C_i a \leq d_i\}$
\emph{ekvivalentan je sljedećem LP-u}:
\[
\min_{x,\ \{\mu_i\}_{i=1,\dots,m}} \ c\trans x
\qquad \text{uz ogr.} \quad
\begin{aligned}
d_i\trans \mu_i &\leq b_i, & & i = 1,\dots,m \\
C_i\trans \mu_i &= x, & & i = 1,\dots,m \\
\mu_i &\geq 0, & & i = 1,\dots,m
\end{aligned}
\]

\textbf{Što ovaj rezultat radi:} unutarnji problem
$\sup_{a \in \mathcal{P}_i} a\trans x$ zamjenjuje njegovim \emph{dualnim}
problemom, pa umjesto maksimizacije unutar ograničenja dobivamo dodatne
varijable $\mu_i$ i dodatna linearna ograničenja. Skup nesigurnosti opisan je
nejednakostima ($C_i a \leq d_i$), a ne vrhovima.
\end{izvor}

\begin{izvor}{Predavanje 8 --- Robusno optimiranje, slajd 4/21 (PDF str.~4)}
\slajd[0.86]{p08-dokaz-dualnost}

\textbf{Transkripcija (dokaz).} Za fiksan $i$, problem
$\sup_{a \in \mathcal{P}_i} a\trans x$ ekvivalentan je problemu
\[
\inf_a \ -a\trans x \quad \text{uz ogr.} \quad C_i a \leq d_i . \qquad (\ast)
\]
Optimizacijski problem $(\ast)$ je LP za koji vrijedi jaka dualnost, pa njegovu
optimalnu vrijednost možemo dobiti iz Lagrangeovog dualnog problema. Za $(\ast)$
imamo Lagrangian
\[
\Lag(a, \mu_i) = -a\trans x + \mu_i\trans (C_i a - d_i),
\]
gdje je $\mu_i \in \R^{m_i}$ vektor dualnih varijabli.

\textbf{Uoči odakle dolazi:} ovo je \emph{isti} aparat kao u Zadatku 1e ---
Lagrangian, dualna funkcija, jaka dualnost. Robusna optimizacija ovdje nije nova
teorija, nego primjena dualnosti na unutarnji problem najgoreg slučaja.
\end{izvor}

\begin{korak}[Zašto su dvije metode ovdje jednake]
\begin{enumerate}[leftmargin=*]
  \item Obje su \textbf{egzaktne} --- nijedna nije aproksimacija. Daju identičan
        dopustivi skup, pa i identično rješenje.
  \item \textbf{Vrhovi} su praktičniji kad je skup nesigurnosti mali pravokutnik:
        $2^2 = 4$ vrha po ograničenju, dakle osam nejednakosti, i problem ostaje
        u istim dvjema varijablama $x_1, x_2$.
  \item \textbf{Dualna reformulacija} je praktičnija kad broj vrhova eksplodira.
        Za nesigurnost u $n$ koeficijenata pravokutnik ima $2^n$ vrhova ---
        eksponencijalno mnogo --- dok dualna metoda dodaje samo $m_i$ novih
        varijabli po ograničenju, dakle linearno mnogo.
  \item Za naš zadatak, s $n = 2$, osam nejednakosti je posve neproblematično, pa
        je odabrana metoda vrhova. Kad bi koeficijenata bilo dvadeset, imali
        bismo preko milijun nejednakosti i metoda bi propala.
\end{enumerate}
\end{korak}

\begin{izgovor}
,,Skup nesigurnosti je pravokutnik, dakle konveksna ljuska svojih četiriju
vrhova. Linearna funkcija po konveksnoj ljusci postiže maksimum u nekom od
generatora, što se dokazuje u jednom retku: zapišem proizvoljnu točku kao
konveksnu kombinaciju vrhova, uvrstim, i iskoristim da su težine nenegativne i
zbrajaju se u jedan. Obrat je trivijalan jer su vrhovi elementi skupa.

Zato je robusno ograničenje egzaktno ekvivalentno s četiri nejednakosti po
ograničenju, ukupno osam. Robusni problem time postaje običan kvadratni program.

Profesor na Predavanju 8 isti rezultat izvodi preko dualnosti unutarnjeg
problema, što daje dodatne dualne varijable umjesto dodatnih nejednakosti. Za
naša dva koeficijenta po retku metoda vrhova je jednostavnija; da ih je bilo
dvadeset, broj vrhova bi bio dva na dvadesetu i morao bih na dualnu formulaciju.``
\end{izgovor}
